
\section{Inference Infra}
\subsection{W8A8 Quantization}
\label{sec:appendix_w8a8}
We adopt the A8W8 SmoothQuant approach~\citep{xiao2023smoothquant}, which leverages a calibration dataset to pre-compute per-channel scaling factors $s$. This enables an equivalent transformation of the form $Y = (X \cdot \text{diag}(s)^{-1}) \cdot (\text{diag}(s)W)$, effectively mitigating the impact of outliers in channel-wise activations.

For calibration, we constructed a dataset encompassing a wide range of usage scenarios, including different task types (\emph{e.g.}, T2V and I2V) and a uniformly sampled step size within the range [12, 32]. Notably, I2V samples constituted approximately 30\% of the dataset. We employed the \texttt{FP8} data type for quantization, as \texttt{INT8} was found to introduce noticeable visual artifacts in the generated videos. Furthermore, we conducted a hyperparameter search over the range $\alpha \in (0.4, 0.6)$, and ultimately selected $\alpha = 0.45$ for SmoothQuant.
All model weights were quantized except for the first and last layers. This quantization strategy led to a 30\% performance improvement without compromising generation quality.


\subsection{Multi-Node Parallel Inference}
\label{sec:appendix_multinode_parallel_inference}
We adopted a multi-node Ulysses-based parallel inference framework across 3 nodes (24 GPUs), where inter-GPU communication and the high computational density of attention emerged as the primary bottlenecks.
Ulysses performs four all-to-all communication steps for the $q$, $k$, $v$, and $o$ tensors. To mitigate communication overhead, we carefully overlapped each communication stage with corresponding computations:

\begin{itemize} 
    \item $v$-communication overlaps with $k$-computation
    \item $k$-communication overlaps with $q$-computation
    \item $q$-communication overlaps with KV cache updates
    \item $o$-communication overlaps with cross-attention computation
\end{itemize}

This overlapping strategy effectively reduced communication overhead to less than 3\% of total execution time.

\begin{figure}[htbp]
    \centering
    \begin{minipage}{1.0\textwidth}
    \begin{subfigure}[b]{1.0\textwidth}
        \centering
        \includegraphics[width=1.0\textwidth]{infra/infer_infra/figures/ulysses_split.pdf}
        \caption{Ulysses context split logic.}
    \end{subfigure}
    \hfill
    \begin{subfigure}[b]{1.0\textwidth}
        \centering
        \includegraphics[width=1.0\textwidth]{infra/infer_infra/figures/cso_split_v2.pdf}
        \caption{CSO context split logic.}
    \end{subfigure}
    \end{minipage}
    \caption{Overview of Ulysses and CSO context split logic.}
    \label{fig:context_split_logic_overview}
\end{figure}

\subsection{Context Shuffle Overlap}
\label{sec:appendix_context_shuffle_overlap}
We optimized based on the Ulysses CP algorithm, striving to overlap communication with computation or data movement. Since the RTX 4090 GPUs communicate via PCIe, which has a relatively low bandwidth, we further proposed the \textbf{CSO (Context Shuffle Overlap)} algorithm to optimize communication more deeply.

The key difference between CSO and Ulysses lies in how the context is partitioned. In Ulysses, all chunks are distributed sequentially across different ranks. In contrast, CSO assigns each rank a partial view of every chunk. As illustrated in Fig.~\ref{fig:context_split_logic_overview}, this alternative partitioning strategy allows CSO to conveniently overlap computation and communication at the chunk level. Assuming we have 5 chunks, the complete process of CSO is as follows:
\begin{itemize} 
    \item $k$-communication and $v$-communication of all chunks overlaps with $q$-computation of all chunks
    \item $q$-communication of chunk 1 overlaps with KV cache updates
    \item $q$-communication of chunk 2 overlaps with $o$-computation of chunk 1
    \item $q$-communication of chunk 3 and $o$-communication of chunk 1 overlaps with $o$-computation of chunk 2
    \item $q$-communication of chunk 4 and $o$-communication of chunk 2 overlaps with $o$-computation of chunk 3
    \item $q$-communication of chunk 5 and $o$-communication of chunk 3 overlaps with $o$-computation of chunk 4
    \item $o$-communication of chunk 4 overlaps with $o$-computation of chunk 5
    \item $o$-communication of chunk 5 overlaps with cross-attention computation
\end{itemize}

In addition, the CSO partition pattern enables communication operations to be split into multiple balanced all-to-all communications. These balanced operations offer better performance compared to unbalanced all-to-all communication and also allow for efficient subsequent merging.
